\documentclass[10pt,a4paper]{amsart}
\usepackage[margin=19mm]{geometry}
\usepackage{amsmath,amssymb,mathtools}
\usepackage[scr=rsfs,cal=euler]{mathalfa}
\usepackage{xfrac}
\usepackage[unicode,hidelinks,bookmarks=false]{hyperref}
\newcommand{\ie}{i.e.\ }
\newcommand{\cf}{cf.\ }
\newcommand{\resp}{respectively\ }
\newcommand{\Subsection}{\subsection}
\providecommand{\nobreakdash}{\nobreak\hbox{-}}
\setlength{\emergencystretch}{2em}
\multlinegap0pt
\usepackage{amssymb}
\newtheorem{lemma}{Lemma}
\newcommand{\C}{{\mathbb C}}
              \newcommand{\E}{{\mathcal E}}
\renewcommand{\P}{{\mathbb P}}
\title{Symplectic configurations: a homological and computer-aided approach}
\author{Weimin Chen}
\date{Source: arXiv:2509.07429v3 (CC BY 4.0)}
\begin{document}
\maketitle

\noindent\emph{Source and licence.} Symplectic configurations: a homological and computer-aided approach. Version arXiv:2509.07429v3 (CC BY 4.0), distributed by arXiv under the Creative Commons Attribution 4.0 International licence, http://creativecommons.org/licenses/by/4.0/, as recorded in the arXiv licence metadata for this exact version and shown on the abstract page https://arxiv.org/abs/2509.07429v3. Full licence notice: \texttt{licenses/arxiv-2509.07429-CC-BY-4.0.txt}.\par\smallskip

\noindent\textit{Editorial scope.} Counted unit: the article's lemma lemma (source0.tex lines 1760--1769). The article's complete original proof of it is delivered with it (source0.tex lines 1771--1829, one \texttt{proof} environment of 59 lines). No mathematics is written, repaired or re-derived: the statement and the proof are the article's own lines, in the article's own notation, and no retained line is substituted or re-flowed.  No further source-local context is needed: every cross-reference of every retained line resolves inside the two delivered spans.  Nothing was omitted from any retained span, so the package carries no omission record.  No retained line contains a citation, so the wrapper carries no bibliography and no bibliography entry.  No retained line contains a figure environment, an image-inclusion command or drawing code, so no image is delivered and the package declares no asset dependency.  Editorial formatting only, in addition: an AMS \texttt{amsart} preamble that restates the article's own declarations and macros used by the retained lines, namely the environments \texttt{lemma} on separate counters, together with the standard packages those lines need.  The retained environments print in this document as Lemma 1 in order of appearance, whereas the article prints the same items with the numbering of its own full document.  The complete native source of the article as distributed in the arXiv e-print for this exact version is bundled unchanged as \texttt{sources/184732/source0.tex}.
\begin{lemma}
Let $(\vec{v}_k^\prime):=R(\gamma)(\vec{v}_k)$, where $\gamma:=H-E_r-E_s-E_t$. 
Then $(\vec{v}_k^\prime)\in\Omega(D)$, i.e., each 
$\vec{v}_k^\prime$ is admissible, if and only if the following conditions hold: for any $k$,
\begin{itemize}
\item [{(i)}] if $a_k=1$, then $b_{kr}+b_{ks}+b_{kt}\leq 2$, and 
\item [{(ii)}] if $a_k=0$, then $b_{kr}+b_{ks}+b_{kt}\leq 0$. 
\end{itemize}
Moreover, the first entry $a_k^\prime$ in each $\vec{v}_k^\prime$ is non-negative. 
\end{lemma}

\begin{proof}
To simplify the notations, we shall drop the index $k$ in $\vec{v}_k$ and $\vec{v}_k^\prime$, simply write it as 
$\vec{v}=(a,b_{1},b_{2}, \cdots, b_{N})$ and $\vec{v}^\prime=(a^\prime,b_{1}^\prime,b_{2}^\prime, \cdots, 
b_{N}^\prime)$.

With this understood, we let $A=aH-\sum_{i=1}^N b_iE_i$ and $A^\prime=a^\prime H-\sum_{i=1}^N b_i^\prime E_i$
where $A^\prime=R(\gamma)(A)=A+(\gamma\cdot A)\gamma$. Note that 
$\gamma\cdot A=a-(b_r+b_s+b_t)$, from which it follows easily that
$$
a^\prime=2a-(b_r+b_s+b_t), \; b_r^\prime=a-(b_s+b_t), \; b_s^\prime=a-(b_r+b_t), \; b_t^\prime=a-(b_r+b_s),
$$
and $b_i^\prime=b_i$ for any $i\neq r,s,t$.

With this understood, consider first the case where $a=1$. Since by assumption $b_{r}+b_{s}+b_{t}\leq 2$, 
it follows easily that $a^\prime\geq 0$, with $a^\prime= 0$ iff $b_{r}+b_{s}+b_{t}=2$. Since $a=1$, it is easy to see
that $b_{r}+b_{s}+b_{t}=2$ if and only if exactly two of $b_{r}, b_{s},b_{t}$ equal to $1$ and the third one equals $0$. It follows easily that $\vec{v}^\prime$ is admissible in this case. 

Next, assume $a=0$. In this case, we have  $b_{r}+b_{s}+b_{t}\leq 0$ by the assumption, which gives 
$a^\prime\geq 0$ immediately,  with $a^\prime= 0$ iff $b_{r}+b_{s}+b_{t}=0$ and $a^\prime>0$  
iff $b_{r}+b_{s}+b_{t}=-1$. It follows easily that $\vec{v}^\prime$ is admissible in this case as well.

Finally, we consider the case where $a\geq 2$. Note that under the condition of $a\geq 2$, one has $a>b_i$ 
for any $i$ (see the proof of Lemma 2.4). The assertion that $\vec{v}^\prime$ is admissible and $a^\prime>0$
follows immediately from the following claim:

\vspace{1mm}

{\bf Claim:} {\it Assume $a\geq 2$. Then for any $i\neq j$, $a\geq b_{i}+b_{j}$ holds true.}

\vspace{1mm}

For a proof, we let $F$ be the component of $D$ whose homological expression is given by the vector 
$\vec{v}$, and let $\hat{F}$ denote the corresponding component in $\hat{D}$, which is a $\hat{J}$-holomorphic curve of degree $a$. 

We begin by recalling the following formula for computing local intersection numbers. Let $C$ be a germ of holomorphic curves (not necessarily irreducible) at $0\in \C^2$, and let $L$ be a germ of embedded holomorphic disk intersecting $C$ only at $0\in \C^2$. We blow up at $0\in \C^2$ and let $E$ be the exceptional $(-1)$-sphere. Let $C^\prime$, $L^\prime$ denote the proper transforms of $C$, $L$ respectively. Then one has
$$
L\cdot C=E\cdot C^\prime+L^\prime\cdot C^\prime. 
$$
In particular, $L\cdot C\geq E\cdot C^\prime$ as $L^\prime\cdot C^\prime\geq 0$. 

Now back to the proof of the claim that $a\geq b_{i}+b_{j}$ for any $i\neq j$. First consider the case where
there exist two distinct minimal classes $E_m$, $E_n$ such that $E_m\leq E_i$, $E_n\leq E_j$.
Let $L$ be the unique degree $1$ $\hat{J}$-holomorphic sphere in $\C\P^2$ passing through 
the points $\hat{E}_m$, $\hat{E}_n$. Then on the one hand, $a=L\cdot \hat{F}$ as $\hat{F}$ is of
degree $a$, and on the other hand, the local intersection number of $L$ with $\hat{F}$ at
$\hat{E}_m$, $\hat{E}_n$ is bounded from below by $b_m, b_n$ respectively. This is because if 
$F^\prime$ denotes the proper transform of $\hat{F}$ after blowing up at $\hat{E}_m$ (resp. 
$\hat{E}_n$), then $b_m=E_m\cdot F^\prime$ (resp. $b_n=E_n\cdot F^\prime$). It follows easily that
$a\geq b_m+b_n\geq b_i+b_j$ by Lemma 3.4(4). 

It remains to consider the case where there is only one minimal class $E_m$ such that $E_m\leq E_i$,
$E_m\leq E_j$. Since $E_i,E_j$ are distinct, there must be a class $E_{i_\alpha}$ which is infinitely near 
to $E_m$ of order $1$ such that $E_{i_\alpha}\leq E_j$ (or $E_i$). With this understood, it suffices to
show that $a\geq b_m+b_{i_\alpha}$. To see this, let $L$ be the unique degree $1$ $\hat{J}$-holomorphic 
sphere in $\C\P^2$ passing through the point $\hat{E}_m$ and tangent to the line $L_\alpha$ in the $4$-ball $B(\hat{E}_m)$ determined by the class $E_{i_\alpha}$ (cf. Description of $\hat{D}$ near the points $\hat{E}_i$, where $E_i\in\E(D)$, i.e., $E_i$ is minimal). Then the local intersection number of $L$ with $\hat{F}$ at the 
point $\hat{E}_m$ equals $b_m+L^\prime\cdot F^\prime$, where $L^\prime,F^\prime$ are the proper transforms of $L$, $\hat{F}$ after blowing up at $\hat{E}_m$. Since $L$ is tangent to $L_\alpha$, $L^\prime$ must be an embedded disk containing 
the center $\hat{E}_{i_\alpha}$ of the $4$-ball $B(\hat{E}_{i_\alpha})$. Consequently, 
$L^\prime\cdot F^\prime\geq b_{i_\alpha}$, and the proof is finished. 
\end{proof}

\end{document}
